% Make \emph meaningful in italic environments (e.g., when using \theoremstyle{plain})
\renewcommand\eminnershape{\itshape\sffamily}

\theoremstyle{plain} % italics
\newtheorem{itheorem}{Theorem}%[section]
\newtheorem{icorollary}{Corollary}%[section]
\newtheorem{theorem}{Theorem}[section]
\newtheorem{lemma}[theorem]{Lemma}
\newtheorem{proposition}[theorem]{Proposition}
\newtheorem{claim}[theorem]{Claim}
\newtheorem{assumption}{Assumption}
\newtheorem{construction}{Construction}
\theoremstyle{definition} % not italics
\newtheorem{idefinition}{Definition}%[section]
\newtheorem{definition}[theorem]{Definition}
\newtheorem{remark}[theorem]{Remark}
\newtheorem{example}[theorem]{Example}
\theoremstyle{remark} %
\newtheorem{case}{Case}
\crefname{assumption}{Assumption}{Assumptions}
\crefname{step}{Step}{Steps}
\crefname{figure}{Figure}{Figures}

\newenvironment{blockquotation}
  {\begin{quote}\begin{singlespace}}
  {\end{singlespace}\end{quote}}
\newcommand{\bq}[1]{\begin{blockquotation}#1\end{blockquotation}}


% Make \emph meaningful in italic environments (e.g., when using \theoremstyle{plain})
\renewcommand\eminnershape{\itshape\sffamily}

\newcommand{\emath}[1]{\ensuremath{#1}}
\newcommand{\newc}[2]{\newcommand{#1}{\ensuremath{#2}\xspace}}
\newcommand{\renewc}[2]{\renewcommand{#1}{\ensuremath{#2}\xspace}}


\renewc{\AA}{\mathbb{A}}
\newc{\CC}{\mathbb{C}}
\newc{\DD}{\mathbb{D}}
\newc{\EE}{\mathbb{E}}
\newc{\FF}{\mathbb{F}}
\newc{\HH}{\mathbb{H}}
\newc{\NN}{\mathbb{N}}
\newc{\PP}{\mathbb{P}}
\newc{\QQ}{\mathbb{Q}}
\newc{\RR}{\mathbb{R}}
\renewc{\SS}{\mathbb{S}}
\newc{\ZZ}{\mathbb{Z}}
\DeclareSymbolFont{bbold}{U}{bbold}{m}{n}
\DeclareSymbolFontAlphabet{\mathbbold}{bbold}
\newc{\One}{\mathbbold{1}}

\newc{\ba}{\bm a}
\newc{\bb}{\bm b}
\newc{\bc}{\bm c}
\newc{\bd}{\bm d}
\newc{\be}{\bm e}
\newc{\bg}{\bm g}
\newc{\bh}{\bm h}
\newc{\bi}{\bm i}
\newc{\bj}{\bm j}
\newc{\bk}{\bm k}
\newc{\bbm}{\bm m}
\newc{\bp}{\bm p}
% \newc{\bq}{\bm q}
\newc{\br}{\bm r}
\newc{\bs}{\bm s}
\newc{\bt}{\bm t}
\newc{\bu}{\bm u}
\newc{\bv}{\bm v}
\newc{\bw}{\bm w}
\newc{\bx}{\bm x}
\newc{\by}{\bm y}
\newc{\bz}{\bm z}

\newc{\bA}{\bm A}
\newc{\bB}{\bm B}
\newc{\bC}{\bm C}
\newc{\bD}{\bm D}
\newc{\bE}{\bm E}
\newc{\bF}{\bm F}
\newc{\bG}{\bm G}
\newc{\bH}{\bm H}
\newc{\bI}{\bm I}
\newc{\bL}{\bm L}
\newc{\bM}{\bm M}
\newc{\bN}{\bm N}
\newc{\bP}{\bm P}
\newc{\bQ}{\bm Q}
\newc{\bR}{\bm R}
\newc{\bS}{\bm S}
\newc{\bT}{\bm T}
\newc{\bU}{\bm U}
\newc{\bV}{\bm V}
\newc{\bW}{\bm W}
\newc{\bX}{\bm X}
\newc{\bY}{\bm Y}
\newc{\bZ}{\bm Z}

\newc{\zero}{\bm{0}}
\newc{\one}{\bm{1}}

\newc{\sA}{\mathcal{A}}
\newc{\sB}{\mathcal{B}}
\newc{\sC}{\mathcal{C}}
\newc{\sD}{\mathcal{D}}
\newc{\sE}{\mathcal{E}}
\newc{\sF}{\mathcal{F}}
\newc{\sG}{\mathcal{G}}
\newc{\sH}{\mathcal{H}}
\newc{\sI}{\mathcal{I}}
\newc{\sJ}{\mathcal{J}}
\newc{\sK}{\mathcal{K}}
\newc{\sL}{\mathcal{L}}
\newc{\sM}{\mathcal{M}}
\newc{\sN}{\mathcal{N}}
\newc{\sO}{\mathcal{O}}
\newc{\sP}{\mathcal{P}}
\newc{\sQ}{\mathcal{Q}}
\newc{\sR}{\mathcal{R}}
\newc{\sS}{\mathcal{S}}
\newc{\sT}{\mathcal{T}}
\newc{\sU}{\mathcal{U}}
\newc{\sV}{\mathcal{V}}
\newc{\sX}{\mathcal{X}}
\newc{\sY}{\mathcal{Y}}

\newc{\fB}{\mathscr{B}}
\newc{\fC}{\mathscr{C}}
\newc{\fE}{\mathscr{E}}
\newc{\fH}{\mathscr{H}}
\newc{\fI}{\mathscr{I}}
\newc{\fP}{\mathscr{P}}
\newc{\fV}{\mathscr{V}}

\DeclareMathOperator{\GOE}{GOE}
\DeclareMathOperator{\sym}{sym}
\DeclareMathOperator{\Part}{Part}
\DeclareMathOperator{\poly}{poly}
\DeclareMathOperator{\Unif}{Unif}
\DeclareMathOperator{\Var}{Var}
\DeclareMathOperator{\Cov}{Cov}
\DeclareMathOperator{\Tr}{Tr}
\DeclareMathOperator{\rank}{rank}
\DeclareMathOperator{\obj}{obj}
\DeclareMathOperator{\diag}{diag}
\DeclareMathOperator{\sgn}{sgn}

\renewc\Re{\operatorname{Re}}
\renewc\Im{\operatorname{Im}}
% These put symbols below.
\newc{\Px}{\mathop{\mathbb{P}}}
\newc{\Ex}{\mathop{\mathbb{E}}}
\newc{\Varx}{\mathop{\mathsf{Var}}}

\newc\numberthis{\addtocounter{equation}{1}\tag{\theequation}}

% allow in-line math breaks
\newc{\CB}{\allowbreak}

% such that in probability statements
\newc{\pST}{\; \middle| \;}
\DeclareMathOperator{\Concat}{\; || \; }
% \DeclareMathOperator{\max}{max}
\newc{\Or}{\vee}
\renewc{\And}{\wedge}
\newc{\Not}{\neg}
\newc{\defeq}{\triangleq}
\newc{\union}{\cup}

% some math macros
\newcommand{\smalllist}[1]{[#1]}
\newcommand{\set}[1]{\ensuremath{\mleft\{#1\mright\}}}
\newcommand{\range}[2]{\set{#1, \dotsc, #2}}
\newcommand{\openopen}[2]{\left[#1, #2\right]}
\newcommand{\openclosed}[2]{\left[#1, #2\right)}
\newcommand{\closedopen}[2]{\left(#1, #2\right]}
\newcommand{\closedclosed}[2]{\left(#1, #2\right)}

\newcommand{\element}[3]{#1_{#2,#3}}

% tuples with in-line math breaks
\newcommand{\tuple}[2]{(#1 ,\CB \dots,\CB #2)}
\newcommand{\pair}[2]{(#1 ,\CB #2)}
\newcommand{\ppair}[2]{\left(#1 ,\CB #2\right)}
\newcommand{\triple}[3]{(#1 ,\CB #2,\CB #3)}
\newcommand{\ptriple}[3]{\left(#1 ,\CB #2,\CB #3\right)}
\newcommand{\quadruple}[4]{(#1 ,\CB #2,\CB #3,\CB #4)}
\newcommand{\pquadruple}[4]{\left(#1 ,\CB #2,\CB #3,\CB #4\right)}
\newcommand{\quintuple}[5]{(#1 ,\CB #2,\CB #3,\CB #4,\CB #5)}
\newcommand{\quintuplenp}[5]{#1 ,\CB #2,\CB #3,\CB #4,\CB #5}
\newcommand{\sextuple}[6]{(#1 ,\CB #2,\CB #3,\CB #4,\CB #5,\CB #6)}
\newcommand{\sextuplenp}[6]{#1 ,\CB #2,\CB #3,\CB #4,\CB #5,\CB #6}
\newcommand{\septuple}[7]{(#1 ,\CB #2,\CB #3,\CB #4,\CB #5,\CB #6,\CB #7)}
\newcommand{\septuplenp}[7]{#1 ,\CB #2,\CB #3,\CB #4,\CB #5,\CB #6,\CB #7}
\newcommand{\octuple}[8]{(#1 ,\CB #2,\CB #3,\CB #4,\CB #5,\CB #6,\CB #7,\CB #8)}
\newcommand{\octuplenp}[8]{#1 ,\CB #2,\CB #3,\CB #4,\CB #5,\CB #6,\CB #7,\CB #8}
\newcommand{\nonuple}[9]{(#1 ,\CB #2,\CB #3,\CB #4,\CB #5,\CB #6,\CB #7,\CB #8,\CB #9)}
\newcommand{\nonuplenp}[9]{#1 ,\CB #2,\CB #3,\CB #4,\CB #5,\CB #6,\CB #7,\CB #8,\CB #9}

% expectation
\newcommand{\prob}[1]{\Pr\left[ #1 \right]}
\newc{\Expectation}{\mathbb{E}}

% random variable
\newc{\randomvariable}{X}

% complexity classes
\newcommand{\class}[1]{\mathsf{#1}}
\newc{\np}{\Class{NP}}
\newcommand{\dtime}[1]{\Class{DTIME}(#1)}

% types
\newc{\bits}{\{0,1\}}
\newc{\strings}{\bits^{*}}
\newc{\ZZq}{\ZZ_{q}}
\newc{\ZZp}{\ZZ_{p}}
\newc{\FFq}{\FF_{q}}
\newc{\FFp}{\FF_{p}}
\newc{\field}{\mathbb{F}}
\newc{\ring}{\sR}
\newc{\ideal}{\mathfrak{n}}
\newc{\idealb}{\mathfrak{b}}
\newc{\idealc}{\mathfrak{c}}
\newc{\primeideal}{\mathfrak{p}}
\newc{\idealnorm}{N}
\newcommand{\content}[1]{\mathsf{\mathsf{cont}\left(#1\right)}}

\DeclarePairedDelimiter{\abs}{\lvert}{\rvert}%
\newcommand{\size}[1]{\abs{#1}}
\newc{\Time}{t}

\newc{\secparam}{\lambda}
\newc{\knowledgeerror}{\kappa}
\renewc{\poly}{\mathrm{poly}}
\newc{\negl}{\mathrm{negl}}
\renewc{\deg}{\mathrm{deg}}
\newc{\view}{\mathrm{View}}

% attackers
\newc{\querysampler}{\mathcal{Q}}
\newc{\adversary}{\mathcal{A}}
\newc{\adversaryb}{\mathcal{B}}
\newc{\adversaryd}{\mathcal{D}}
\newc{\experiment}{\mathsf{Exp}}
\newc{\simulator}{\mathcal{S}}
\newc{\real}{\mathsf{Real}}
\newc{\extractor}{\mathcal{E}}
\newc{\extractoradversary}{\mathcal{E}_{\adversary}}
\newc{\challenger}{\mathcal{C}}

\newc{\prover}{\mathsf{P}}
\newc{\verifier}{\mathsf{V}}

\makeatletter
\long\def\defcommand#1{%
  \@ifnextchar[{\defcommand@i{#1}}{\defcommand@ii{#1}}%
}
\long\def\defcommand@i#1[#2]#3{%
  \ifdefined#1%
    \renewcommand{#1}[#2]{#3}%
  \else%
    \newcommand{#1}[#2]{#3}%
  \fi%
}
\long\def\defcommand@ii#1#2{%
  \ifdefined#1%
    \renewcommand{#1}{#2}%
  \else%
    \newcommand{#1}{#2}%
  \fi%
}
\makeatother

\newcommand{\defc}[2]{%
  \ifdefined#1%
    \renewc{#1}{#2}%
  \else%
    \newc{#1}{#2}%
  \fi%
}

\newcommand{\norm}[1]{\ensuremath{\left\lVert #1 \right\rVert}}
\newcommand{\ceil}[1]{\ensuremath{\left\lceil #1 \right\rceil}}

\makeatletter

% Add a period at the end unless the string already ends with punctuation.
% (based on code by Uwe Lueck; see dtrt.sty)
\def\dt@TestPunct#1{\ifx\dt@TestPunct#1\dt@TestPunct\else\dt@Puncttrue\fi}
\def\dt@TestDot#1.\dt@TestEnd#2\dt@TestStop{\dt@TestPunct{#2}}
\def\dt@TestQuestM#1?\dt@TestEnd#2\dt@TestStop{\dt@TestPunct{#2}}
\def\dt@TestExclaM#1!\dt@TestEnd#2\dt@TestStop{\dt@TestPunct{#2}}
\def\dt@TestColonM#1:\dt@TestEnd#2\dt@TestStop{\dt@TestPunct{#2}}
\newif\ifdt@Punct
\def\dt@MaybeAddPunct#1{%
  \dt@Punctfalse%
  \dt@TestDot#1\dt@TestEnd.\dt@TestEnd\dt@TestStop%
  \dt@TestQuestM#1\dt@TestEnd?\dt@TestEnd\dt@TestStop%
  \dt@TestExclaM#1\dt@TestEnd!\dt@TestEnd\dt@TestStop%
  \dt@TestColonM#1\dt@TestEnd:\dt@TestEnd\dt@TestStop%
  #1\ifdt@Punct \else .\fi%
}

% Macros for ignoring spaces and paragraphs.
% The first ignores all spaces and implicit paragraphs (empty separation line),
% the second also explicit \par.
% (Source noted in dtrt.sty)
\def\dt@ignorespacesandimplicitepars{%
  \begingroup \catcode13=10 \@ifnextchar\relax {\endgroup}%
  {\endgroup}%
}
\def\dt@ignorespacesandallpars{%
  \begingroup \catcode13=10 \@ifnextchar\par
  {\endgroup\expandafter\ignorespacesandallpars\@gobble}%
  {\endgroup}%
}

% \parheadnoperiod{text}: parhead without extra punctuation.
\newcommand{\parheadnoperiod}[1]{%
  \smallskip \noindent
  {\bfseries\boldmath\ignorespaces {#1}}%
  \hskip 0.9em plus 0.3em minus 0.3em%
  \dt@ignorespacesandimplicitepars%
}

% \parhead{text}: appends period if needed.
\newcommand{\parhead}[1]{%
  \parheadnoperiod{\dt@MaybeAddPunct{#1}}%
  \dt@ignorespacesandimplicitepars%
}

\makeatother

% Covers:  array, matrix, bmatrix, pmatrix, vmatrix, Vmatrix, cases,    tabular, tabular*, longtable
\renewcommand{\arraystretch}{0.60}

\BeforeBeginEnvironment{table}{\begin{singlespace}\vspace*{-\baselineskip}}
\AfterEndEnvironment{table}{\end{singlespace}\noindent\ignorespaces}
\BeforeBeginEnvironment{longtable}{\begin{singlespace}\vspace*{-\baselineskip}}
\AfterEndEnvironment{longtable}{\end{singlespace}\noindent\ignorespaces}

\BeforeBeginEnvironment{align}{\begin{singlespace}\vspace*{-\baselineskip}}
\AfterEndEnvironment{align}{\end{singlespace}\noindent\ignorespaces}
\BeforeBeginEnvironment{align*}{\begin{singlespace}\vspace*{-\baselineskip}}
\AfterEndEnvironment{align*}{\end{singlespace}\noindent\ignorespaces}

\BeforeBeginEnvironment{alignat}{\begin{singlespace}\vspace*{-\baselineskip}}
\AfterEndEnvironment{alignat}{\end{singlespace}\noindent\ignorespaces}
\BeforeBeginEnvironment{alignat*}{\begin{singlespace}\vspace*{-\baselineskip}}
\AfterEndEnvironment{alignat*}{\end{singlespace}\noindent\ignorespaces}

\BeforeBeginEnvironment{flalign}{\begin{singlespace}\vspace*{-\baselineskip}}
\AfterEndEnvironment{flalign}{\end{singlespace}\noindent\ignorespaces}
\BeforeBeginEnvironment{flalign*}{\begin{singlespace}\vspace*{-\baselineskip}}
\AfterEndEnvironment{flalign*}{\end{singlespace}\noindent\ignorespaces}

\BeforeBeginEnvironment{gather}{\begin{singlespace}\vspace*{-\baselineskip}}
\AfterEndEnvironment{gather}{\end{singlespace}\noindent\ignorespaces}
\BeforeBeginEnvironment{gather*}{\begin{singlespace}\vspace*{-\baselineskip}}
\AfterEndEnvironment{gather*}{\end{singlespace}\noindent\ignorespaces}

\BeforeBeginEnvironment{multline}{\begin{singlespace}\vspace*{-\baselineskip}}
\AfterEndEnvironment{multline}{\end{singlespace}\noindent\ignorespaces}
\BeforeBeginEnvironment{multline*}{\begin{singlespace}\vspace*{-\baselineskip}}
\AfterEndEnvironment{multline*}{\end{singlespace}\noindent\ignorespaces}

% Optional, only if you still use eqnarray:
\BeforeBeginEnvironment{eqnarray}{\begin{singlespace}\vspace*{-\baselineskip}}
\AfterEndEnvironment{eqnarray}{\end{singlespace}\noindent\ignorespaces}
\BeforeBeginEnvironment{eqnarray*}{\begin{singlespace}\vspace*{-\baselineskip}}
\AfterEndEnvironment{eqnarray*}{\end{singlespace}\noindent\ignorespaces}

% Wrap any section/block in this to force single spacing inside it
\newenvironment{singlespacedblock}
  {\begingroup\singlespacing}
  {\endgroup}